\documentclass[10pt,a4paper]{amsart}
\usepackage[margin=19mm]{geometry}
\usepackage{amsmath,amssymb,mathtools}
\usepackage[scr=rsfs,cal=euler]{mathalfa}
\usepackage{xfrac}
\usepackage[unicode,hidelinks,bookmarks=false]{hyperref}
\newcommand{\ie}{i.e.\ }
\newcommand{\cf}{cf.\ }
\newcommand{\resp}{respectively\ }
\newcommand{\Subsection}{\subsection}
\providecommand{\nobreakdash}{\nobreak\hbox{-}}
\setlength{\emergencystretch}{2em}
\multlinegap0pt
\usepackage[shortlabels]{enumitem}
\usepackage{amssymb}
\usepackage{bbm}
\usepackage{xcolor}
\usepackage{float}
\usepackage{tikz}
\usepackage{booktabs}
\usepackage{esint}
\usepackage{multirow}
\newtheorem{theo}{Theorem}
\newtheorem{lemma}{Lemma}
\newcommand{\Teps}{{G_\varepsilon}}
\newcommand{\dXY}{{\mathrm{dist}(X,Y)}}
\newcommand{\jnabla}{{\langle\nabla\rangle}}
\title{Exponential Decay outside of the Light Cone for the Pseudo-Relativistic Non-Autonomous Schrödinger Equation}
\author{S\'ebastien Breteaux, J\'er\'emy Faupin, Viviana Grasselli}
\date{Source: arXiv:2512.20488v1 (CC BY 4.0)}
\begin{document}
\maketitle

\noindent\emph{Source and licence.} Exponential Decay outside of the Light Cone for the Pseudo-Relativistic Non-Autonomous Schrödinger Equation. Version arXiv:2512.20488v1 (CC BY 4.0), distributed by arXiv under the Creative Commons Attribution 4.0 International licence, http://creativecommons.org/licenses/by/4.0/, as recorded in the arXiv licence metadata for this exact version and shown on the abstract page https://arxiv.org/abs/2512.20488v1. Full licence notice: \texttt{licenses/arxiv-2512.20488-CC-BY-4.0.txt}.\par\smallskip

\noindent\textit{Editorial scope.} Counted unit: the article's lemma lm:bound\_domain (source0.tex lines 596--605). The article's complete original proof of it is delivered with it (source0.tex lines 607--679, one \texttt{proof} environment of 73 lines). No mathematics is written, repaired or re-derived: the statement and the proof are the article's own lines, in the article's own notation, and no retained line is substituted or re-flowed.  Retained uncounted as the source-local context the proof needs: lines 167--174; lines 437--442.  Nothing was omitted from any retained span, so the package carries no omission record.  No retained line contains a citation, so the wrapper carries no bibliography and no bibliography entry.  No retained line contains a figure environment, an image-inclusion command or drawing code, so no image is delivered and the package declares no asset dependency.  Editorial formatting only, in addition: an AMS \texttt{amsart} preamble that restates the article's own declarations and macros used by the retained lines, namely the environments \texttt{theo}, \texttt{lemma} on separate counters, together with the standard packages those lines need.  The retained environments print in this document as Theo 1, Lemma 1 in order of appearance, whereas the article prints the same items with the numbering of its own full document.  The complete native source of the article as distributed in the arXiv e-print for this exact version is bundled unchanged as \texttt{sources/191571/source0.tex}.
Let $\|\cdot\|_{L^2}$ be the operator norm on $L^2(\mathbb{R}^d)$ and let $\mathrm{dist}(X,Y)$ be the distance between two subsets $X$ and $Y$ of $\mathbb{R}^d$.  Our main result provides, for any $t$ in $[0,T]$, an estimate on the norm $\|\mathbf{1}_YU_t\mathbf{1}_X\|_{\mathcal{B}(L^2)}$ of the unitary propagator generated by $(\jnabla+V_t)_{t\in[0,T]}$ composed with the characteristic functions $\mathbf{1}_X$ and $\mathbf{1}_Y$ of any convex subsets~$X$ and $Y$ of $\mathbb{R}^d$:
%
\begin{theo}\label{th:max-vel}
Let $T>0$ and $(V_t)_{t\in[0,T]}$ be a family of real-valued potentials such that the family~$(U_t)_{t\in[0,T]}$ is a propagator generated by $(\langle\nabla\rangle+V_t)_{t\in[0,T]}$. If $X$ and $Y$ are convex subsets of~$\mathbb{R}^d$, then
\begin{equation}\label{eq:main-est}
        \forall t\in [0,T]\,,\quad \|\mathbf{1}_Y U_t \mathbf{1}_X\|_{\mathcal{B}(L^2)} \le e^{t-\dXY}.
    \end{equation}
\end{theo}

	\begin{lemma}\label{lm:unif-bound}
	    For all $\varepsilon>0$, $\Teps$ extends to a bounded operator on $L^2$, with
	    \begin{equation*}
	        \sup_{\varepsilon>0}\|\Teps\|_{\mathcal{B}(L^2)}<\infty.
	    \end{equation*}
	\end{lemma}

	\begin{lemma}\label{lm:bound_domain}Under the assumptions of Theorem~\ref{th:max-vel}, for all $t$ in $[0,T]$, we have
    \begin{equation*}
        \mathrm{Ran}(U_te^{\ell(x)})\subset\mathcal{D}(e^{-\ell(x)}),
    \end{equation*}
    and $e^{-\ell(x)}U_te^{\ell(x)}$ extends to a bounded operator on $L^2$. Moreover,
    \begin{equation*}
        e^{-\ell_\varepsilon(x)}U_te^{\ell_\varepsilon(x)} \to e^{-\ell(x)}U_te^{\ell(x)} \quad \text{as}\quad \varepsilon\to0,
    \end{equation*}
    strongly in $\mathcal{B}(L^2)$. 
\end{lemma}

	\begin{proof}
	    Let $\varphi$ in $\mathcal{C}^\infty_0$ and  $\varepsilon>0$.  We  compute
	    \begin{align*}
    &\big\|e^{-\ell_\varepsilon(x)}U_te^{\ell_\varepsilon(x)}\varphi\big\|_{L^2}^2\\
    &=\|\varphi\|_{L^2}^2+2\mathrm{Re}\int_0^t\big\langle e^{-\ell_\varepsilon(x)}U_\tau e^{\ell_\varepsilon(x)}\varphi \, , \, e^{-\ell_\varepsilon(x)} (-iH_\tau)U_\tau e^{\ell_\varepsilon(x)}\varphi\big\rangle\mathrm{d}\tau\\
    &=\|\varphi\|_{L^2}^2+2\mathrm{Re}\int_0^t\big\langle e^{-\ell_\varepsilon(x)}U_\tau e^{\ell_\varepsilon(x)}\varphi \, , \, e^{-\ell_\varepsilon(x)} (-i\langle\nabla\rangle)e^{\ell_\varepsilon(x)}e^{-\ell_\varepsilon(x)}U_\tau e^{\ell_\varepsilon(x)}\varphi\big\rangle\mathrm{d}\tau,
\end{align*}
where in the second equality we used that
\begin{equation*}
    \mathrm{Re}\big\langle e^{-\ell_\varepsilon(x)}U_\tau e^{\ell_\varepsilon(x)}\varphi \, , \, e^{-\ell_\varepsilon(x)} (-iV_t)e^{\ell_\varepsilon(x)}e^{-\ell_\varepsilon(x)}U_\tau e^{\ell_\varepsilon(x)}\varphi\big\rangle=0.
\end{equation*}
Since
\begin{align*}
    2\mathrm{Re}\, e^{-\ell_\varepsilon(x)} (-i\langle\nabla\rangle)e^{\ell_\varepsilon(x)}&=2\mathrm{Im}\,\big( e^{-\ell_\varepsilon(x)}\langle\nabla\rangle e^{\ell_\varepsilon(x)}\big)=2\Teps, 
\end{align*}
we can rewrite the equality above as
	\begin{align}\label{eq:Gronwall}
    &\big\|e^{-\ell_\varepsilon(x)}U_te^{\ell_\varepsilon(x)}\varphi\big\|_{L^2}^2=\|\varphi\|_{L^2}^2+2\int_0^t\big\langle e^{-\ell_\varepsilon(x)}U_\tau e^{\ell_\varepsilon(x)}\varphi \, , \, \Teps e^{-\ell_\varepsilon(x)}U_\tau e^{\ell_\varepsilon(x)}\varphi\big\rangle\mathrm{d}\tau.
\end{align}
Now we can use Lemma \ref{lm:unif-bound} to deduce that
\begin{align*}
    &\big\|e^{-\ell_\varepsilon(x)}U_te^{\ell_\varepsilon(x)}\varphi\big\|_{L^2}^2
    \le\|\varphi\|_{L^2}^2+C\int_0^t\big\| e^{-\ell_\varepsilon(x)}U_\tau e^{\ell_\varepsilon(x)}\varphi\big\|_{L^2}^2\mathrm{d}\tau,
\end{align*}
for some positive constant $C$ independent of $\varepsilon$. Hence, by Gronwall's inequality,
\begin{align}\label{eq:Gronwall-eps}
    &\big\|e^{-\ell_\varepsilon(x)}U_te^{\ell_\varepsilon(x)}\varphi\big\|_{L^2}^2
    \le e^{Ct}\|\varphi\|_{L^2}^2,
\end{align}
uniformly in $\varepsilon$. By density, this inequality can be extended to any $\varphi$ in $L^2$.

For all $\varphi$ in $\mathcal{D}(e^{\ell(x)})$ and $\psi$ in $\mathcal{D}(e^{-\ell(x)})$, using that $\|U_t\|_{\mathcal{B}(L^2)}=1$, we can then write
	    \begin{align*}
	        \big|\langle e^{-\ell(x)}\psi , U_te^{\ell(x)}\varphi\rangle\big|
	        &=\lim_{\varepsilon\to0}\big|\langle \psi , e^{-\ell_\varepsilon(x)}U_te^{\ell_\varepsilon(x)}\varphi\rangle\big|
	        \le e^{\frac12Ct}\|\psi\|_{L^2}\|\varphi\|_{L^2}.
	   \end{align*}
This proves that $\mathrm{Ran}(U_te^{\ell(x)})\subset\mathcal{D}(e^{-\ell(x)})$ and that $e^{-\ell(x)}U_te^{\ell(x)}$ extends to a bounded operator on $L^2$.

To prove the strong convergence, consider now $\varphi$ in $L^2$. Let $\delta>0$ and let $\varphi_\delta$ in $\mathcal{C}_0^\infty$ be such that $\|\varphi-\varphi_\delta\|_{L^2}\le\delta$. We write
\begin{align}
   & e^{-\ell(x)}U_te^{\ell(x)}\varphi-e^{-\ell_\varepsilon(x)}U_te^{\ell_\varepsilon(x)}\varphi\notag \\
    & = (e^{-\ell(x)}-e^{-\ell_\varepsilon(x)})U_te^{\ell(x)}\varphi_\delta-e^{-\ell_\varepsilon(x)}U_te^{\ell(x)}(e^{\ell_\varepsilon(x)}e^{-\ell(x)}-1)\varphi_\delta\notag \\
    &\quad+ \big(e^{-\ell(x)}U_te^{\ell(x)}-e^{-\ell_\varepsilon(x)}U_te^{\ell_\varepsilon(x)}\big)(\varphi-\varphi_\delta),\label{eq:strongconv1}
\end{align}
and estimate the $L^2$-norm of each term on the right-hand side separately. Using \eqref{eq:Gronwall-eps}, the third term is bounded by
\begin{align}
   \big\| \big(e^{-\ell(x)}U_te^{\ell(x)}-e^{-\ell_\varepsilon(x)}U_te^{\ell_\varepsilon(x)}\big)(\varphi-\varphi_\delta)\big\|_{L^2}\le C_t\delta, \label{eq:strongconv2}
\end{align}
uniformly in $\varepsilon>0$. To estimate the second term, observing that $-\ell_\varepsilon(x)\le\max(-\ell(x),0)$, we can bound
\begin{align*}
   & \big\|e^{-\ell_\varepsilon(x)}U_te^{\ell(x)}(e^{\ell_\varepsilon(x)}e^{-\ell(x)}-1)\varphi_\delta\big\|_{L^2}\\
   &\le\big\|e^{-\ell(x)}U_te^{\ell(x)}(e^{\ell_\varepsilon(x)}e^{-\ell(x)}-1)\varphi_\delta\big\|_{L^2}+\big\|U_te^{\ell(x)}(e^{\ell_\varepsilon(x)}e^{-\ell(x)}-1)\varphi_\delta\big\|_{L^2}\\
   &\le C\big\|(e^{\ell_\varepsilon(x)}e^{-\ell(x)}-1)\varphi_\delta\big\|_{L^2}+\big\|(e^{\ell_\varepsilon(x)}-e^{\ell(x)})\varphi_\delta\big\|_{L^2},
\end{align*}
for some positive constant $C$, where we used that $e^{-\ell(x)}U_te^{\ell(x)}$ belongs to $\mathcal{B}(L^2)$ and $U_t$ is unitary in the second line. Since $\varphi_\delta$ is compactly supported, we deduce from Lebesgue's dominated convergence Theorem that
\begin{align}
   & \big\|e^{-\ell_\varepsilon(x)}U_te^{\ell(x)}(e^{\ell_\varepsilon(x)}e^{-\ell(x)}-1)\varphi_\delta\big\|_{L^2}\to 0 \quad\text{as}\quad  \varepsilon\to0.\label{eq:strongconv3}
\end{align}
It remains to estimate the first term on the right-hand side of \eqref{eq:strongconv1}. Since (using the bound $-\ell_\varepsilon(x)\le\max(-\ell(x),0)$) we have
\begin{align*}
    &\big|(e^{-\ell(x)}-e^{-\ell_\varepsilon(x)})(U_te^{\ell(x)}\varphi_\delta)(x)\big|\le2e^{-\ell(x)}|U_te^{\ell(x)}\varphi_\delta|(x)+|U_te^{\ell(x)}\varphi_\delta|(x),
\end{align*}
we can apply again Lebesgue's dominated convergence Theorem to deduce that
\begin{align}
    \big\|(e^{-\ell(x)}-e^{-\ell_\varepsilon(x)})U_te^{\ell(x)}\varphi_\delta\big\|_{L^2}\to0 \quad\text{as}\quad \varepsilon\to0.\label{eq:strongconv4}
\end{align}
Putting together \eqref{eq:strongconv1}--\eqref{eq:strongconv4}, we have shown that
\begin{align*}
   \limsup_{\varepsilon\to0} \big\|e^{-\ell(x)}U_te^{\ell(x)}\varphi-e^{-\ell_\varepsilon(x)}U_te^{\ell_\varepsilon(x)}\varphi\big\|_{L^2}\le C\delta.
\end{align*}
Since $\delta>0$ is arbitrary, this concludes the proof.
	\end{proof}

\end{document}
